
\subsubsection{Preference Optimization}

Direct Preference Optimization \citep{Rafailov2023DPO} eliminates the instability and complexity of reward model training by directly optimizing language models on preference pairs, achieving 57.5\% win rate versus supervised fine-tuning baselines on dialogue tasks. DPO reparameterizes the RLHF objective to optimize log probability ratios between chosen and rejected responses, implicitly learning a reward model without explicit reward training. However, DPO provides only AI-to-Human alignment: once trained, the model's behavior is fixed to the learned preference distribution with no mechanism for user-specific customization.

Related methods including PPO-based RLHF \citep{Ouyang2022InstructGPT} and Constitutional AI \citep{Bai2022Constitutional} share this limitation. These approaches optimize toward aggregate human values but cannot accommodate diverse individual preferences post-training without complete retraining.

\subsubsection{Attribute-Conditioned Generation}

SteerLM \citep{Dong2023SteerLM} enables users to steer model outputs along interpretable dimensions through attribute-conditioned supervised fine-tuning, achieving 87\% steering accuracy with minimal latency overhead. By conditioning generation on user-specified attribute levels during inference, SteerLM provides Human-to-AI alignment through explicit user control. However, SteerLM operates independently of preference optimization---there is no guarantee that steerable outputs satisfy quality constraints learned from preference data.

Controllable text generation methods more broadly (e.g., CTRL, PPLM, GeDi) enable steering along predefined attributes but typically lack integration with preference-based alignment. Length-normalized DPO \citep{Park2024LengthNormalized} separates length from quality but targets only a single attribute dimension without extending to multi-attribute user control.

\subsubsection{Multi-Task Learning}

Multi-task learning theory provides the foundation for our gradient compatibility analysis. Nash-MTL \citep{Navon2022NashMTL} formulates multi-task optimization as a bargaining game, establishing that when task gradients have positive cosine similarity (angles less than 90°), joint optimization can achieve Pareto improvements. PCGrad \citep{Yu2020PCGrad} and Gradient Surgery \citep{Wang2020GradientSurgery} address catastrophic interference by projecting conflicting gradients when angles exceed 120°, the established threshold for destructive task conflict.

Our contribution demonstrates that DPO and attribute objectives do not require such intervention---their natural gradient alignment (78.5° mean angle) enables joint optimization with simple weighted summation. Representation Surgery \citep{Yang2024RepresentationSurgery} shows that multi-task models can maintain task-specific representations when tasks share complementary structure. Our linear probing analysis achieving 100\% preference classification accuracy from joint model hidden states supports this finding.

\subsubsection{Research Gap}

No previous work has demonstrated that preference optimization and attribute conditioning can be jointly optimized in a single training run, nor quantified their gradient compatibility to predict multi-task feasibility. Sequential training introduces catastrophic forgetting risks and computational overhead, while standalone approaches force users to choose between quality guarantees and customization control. Our work fills this gap through gradient-level analysis providing quantitative evidence that DPO's implicit reward modeling and attribute conditioning's explicit user control are mathematically compatible objectives.
